\chapter{MCMC Family I: Metropolis and Adaptive}
\label{ch:mcmc-1}

\newthought{Markov chain Monte Carlo} (\textsc{mcmc}) draws samples from the posterior by taking a
guided random walk: from the current point it proposes a nearby point and accepts or rejects it so
that, over time, the chain visits each region in proportion to its likelihood. This chapter covers
the baseline random-walk method and three adaptive variants that tune the proposal for you.

\section{The shared format}
\label{sec:mcmc1-format}

Every sampler in this chapter (and in Chapters~\ref{ch:mcmc-2}--\ref{ch:mcmc-3}) reads the same
three \yamlkey{Bounds} keys:

\begin{description}[style=nextline, leftmargin=1.2em]
  \item[\yamlkey{num\_chains}] how many chains to run in parallel ($\geq 1$). Multiple chains
        improve coverage and let you check convergence between them.
  \item[\yamlkey{num\_iters}] total iterations ($\geq 1$).
  \item[\yamlkey{proposal\_scale}] the size of each proposed step ($>0$).
\end{description}

\begin{yamlwin}
Sampling:
  Method: "MCMC"
  Variables: [ ... ]
  Bounds:
    num_chains: 8
    num_iters: 600
    proposal_scale: 0.06
  LogLikelihood:
    - {name: "LogL", expression: "..."}
\end{yamlwin}

\begin{jtip}
\yamlkey{proposal\_scale} is the dial that matters most. Too large and almost every step is
rejected; too small and the chain crawls. Aim for an acceptance rate around 20--40\%; the adaptive
methods below tune this for you.
\end{jtip}

\section{MCMC (random-walk Metropolis)}
\label{sec:s-mcmc}

\paragraph{How it works} The baseline method. From the current point it proposes a Gaussian step of
size \yamlkey{proposal\_scale}, computes the likelihood there, and accepts the move with the
Metropolis probability $\min(1, L_\text{new}/L_\text{old})$---always stepping uphill, sometimes
downhill. The result is a chain whose density traces the posterior.

\paragraph{When to use} A simple, robust starting point for low-to-medium dimensional posteriors.

\begin{jnote}
\alias{ToyMCMC}{MCMC} selects this same sampler under a lightweight name, handy for tests.
\end{jnote}

\section{AMMCMC (adaptive Metropolis)}
\label{sec:s-ammcmc}

\paragraph{How it works} Plain random-walk \textsc{mcmc} struggles when parameters are correlated,
because a single fixed step size cannot match a tilted, stretched posterior. \sampler{AMMCMC}
\emph{learns the proposal covariance} from the chain's own history: as it accumulates samples it
shapes future proposals to follow the posterior's correlations, so it explores far more
efficiently without hand-tuning.

\paragraph{When to use} Posteriors with strong parameter correlations, or when tuning
\yamlkey{proposal\_scale} by hand is painful.

\section{RobustAM (robust adaptive Metropolis)}
\label{sec:s-robustam}

\paragraph{How it works} \sampler{RobustAM} adapts the proposal to hit a \emph{target acceptance
rate}: after each step it nudges the proposal scale/shape up or down so the chain settles at an
efficient acceptance fraction. This makes it forgiving of a poor initial \yamlkey{proposal\_scale}.

\paragraph{When to use} When you want adaptation that is stable from a bad starting guess, with
less sensitivity to the initial scale than plain adaptive Metropolis.

\section{DRAM (delayed rejection + adaptive Metropolis)}
\label{sec:s-dram}

\paragraph{How it works} \sampler{DRAM} combines two ideas. \emph{Delayed rejection}: when a
proposed step is rejected, instead of staying put the chain immediately tries a second, usually
smaller, proposal---and possibly a third---before giving up, which rescues moves in tight regions.
\emph{Adaptive Metropolis}: the proposal covariance is learned from history, as in
\sampler{AMMCMC}. Together they explore awkward posteriors efficiently.

\paragraph{When to use} Difficult posteriors where plain Metropolis gets stuck and you want both
better local acceptance and automatic tuning.

\section{Summary}
\label{sec:mcmc1-summary}

All four share \yamlkey{num\_chains}/\yamlkey{num\_iters}/\yamlkey{proposal\_scale}. Start with
\sampler{MCMC}; switch to \sampler{AMMCMC} or \sampler{DRAM} for correlated or awkward posteriors,
and \sampler{RobustAM} when you want acceptance-rate-targeted adaptation. After any run, check
trace plots, acceptance rate, and between-chain agreement.
